\documentclass[10pt]{article}
\usepackage[margin=1in]{geometry}
\usepackage{amsmath,amssymb,booktabs,hyperref}
\title{Read-Side Cross-Layer KV Sharing for Prefix-Hit Re-Prefill:\\Recovery Benefits and Quality Trade-offs}
\author{Author information omitted}
\date{Working draft based on the v2.6 evidence; not a submission template}
\begin{document}
\maketitle
\begin{abstract}
Prefix cache hits avoid recomputing shared context, but restoring layer-specific
key--value states can still dominate time to first token. We study a frozen-backbone
approach that moves cross-layer compensation to the query and attention-output
paths, allowing multiple logical layers to read a shared physical prefix cache
without materializing complete consumer-specific caches. A bounded subset of the
originally selected prefix positions can remain native to each consumer. Current
experiments on Qwen2.5-7B show lower observed SSD-cold re-prefill latency in all
16 task--budget cells, with reductions of 17.87--42.01\%. However, accuracy declines
in 15 of these cells. A separately blocked TREC/25\% study shows 39.15\% lower cold
latency under matched cache-admission policies. These results support a recovery
cost--quality trade-off rather than universal lossless compression. External
cross-layer baselines and quality-matched operating points remain to be completed.
\end{abstract}

\section{Problem and scope}
We consider requests whose historical prefix has already been computed and
persisted. Online work includes selecting and restoring prefix KV, processing the
new query, and producing the first-token logits. We do not claim savings in prefix
projection work already avoided by the cache hit. Current query and generated-token
KV remain layer-private. All decoder layers, residual paths, and feed-forward
networks execute normally. This scope follows the re-prefill setting distinguished
by ContiguousKV~\cite{contiguouskv}.

\section{Method}
Let $a(\ell)$ denote the physical source for logical consumer layer $\ell$.
Within a KV head, calibrated matrices and offsets define the affine approximation
\begin{align}
 \widehat K_\ell &= K_{a(\ell)}A_\ell+\mathbf{1}b^K_\ell,\\
 \widehat V_\ell &= V_{a(\ell)}B_\ell+\mathbf{1}b^V_\ell.
\end{align}
The approximation is evaluated without reconstructing these complete tensors:
$Q_\ell\widehat K_\ell^T=(Q_\ell A_\ell^T)K_{a(\ell)}^T+
(Q_\ell(b^K_\ell)^T)\mathbf{1}^T$. Likewise, a shared-value weighted aggregate is
transformed after aggregation. These algebraic identities are not themselves a
novelty claim; approximation quality must be measured separately.

The originally selected support $\mathcal I_\ell$ is partitioned into a shared
part $\mathcal C_\ell$ and an optional native subset $\mathcal R_\ell$, with
$\mathcal I_\ell=\mathcal C_\ell\dot\cup\mathcal R_\ell$. No extra tokens are
added. The current implementation uses the last four \emph{selected physical
blocks}, not necessarily an absolute-position sliding window.
Let $u_C,z_C$ and $u_N,z_N$ be conditional outputs and log-normalizers for the
shared and native/private segments. The shared key offset is included in $z_C$.
With $\beta=\exp(z_C-\operatorname{logaddexp}(z_C,z_N))$, the result is
\begin{equation}
 O_\ell=\beta(u_CB_\ell+b^V_\ell)+(1-\beta)u_N.
\end{equation}
This preserves joint normalization between historical and new information.

\section{Offline calibration and online ownership}
The backbone is frozen; only read-side parameters are calibrated offline.
Source maps and readers are currently specific to the model, prefix, sparse
budget, and calibration protocol. Online P8 selection still depends on the current
query. Multiple consumers share immutable source storage; native subsets are loaded
and admitted separately; the private tail never becomes shared across layers.
The durable original 28-layer cache is retained in these experiments, so request
buffer savings are not persistent-storage reductions.

\section{Evidence and limitations}
The v2.6 matrix spans SST-2, SUBJ, TREC, and RTE at 5, 10, 25, and 50\% token
budgets. Accuracy uses the complete declared evaluation partitions; matrix latency
uses 128 paired requests per cell and does not certify each cell independently.
The matched executor control uses the project's optimized P8 path, not automatically
a faithful full ContiguousKV implementation.

\begin{table}[h]
\centering
\begin{tabular}{lrr}
\toprule
Frozen TREC/25\% control & Native identity & Shared reader\\
\midrule
Accuracy (500 historical queries) & 68.60\% & 73.00\%\\
Cold TTFT, matched reuse admission (ms) & 1745.57 & 1062.24\\
Ordinary-cache TTFT (ms) & 261.93 & 201.68\\
\bottomrule
\end{tabular}
\caption{Eight-block primary study and historical quality case. This special case
does not override the 15/16 quality declines in the broader matrix. Source: the
user-provided v2.6 experiment report.}
\end{table}

A fixed-checkpoint unit-reader intervention reaches 16.80\%, showing the importance
of the learned reader as a whole. Replacing native patch contents with affine
source predictions reaches 76.80\% in this one case, higher than native patches.
Thus the data do not establish that native content is indispensable. In exact
arithmetic, the proxy construction and frozen-reader all-shared attention represent
the same affine memory; finite-precision operation order and selector feedback
remain separate effects to test.

At 5\% budget, reader and packed-storage overhead produces negative net subsystem
memory savings in three tasks. The broad matrix must therefore report total
subsystem storage, quality, and latency together. No claim is made of universal
non-inferiority or a larger-model result.

\section{Required comparisons}
ContiguousKV addresses granular recovery and prefetching~\cite{contiguouskv};
KVSharer searches a direct layer-sharing strategy~\cite{kvsharer}; xKV extracts
cross-layer token bases and includes selective reconstruction~\cite{xkv}.
A fixed unit-reader ablation is not a substitute for KVSharer's search. Our current
CLF quality reference is not the full xKV-SR serving engine. Before a submission,
compare these methods at matched quality or total resource budgets in a clearly
specified prefix-hit adaptation, preserving their original algorithmic choices
where possible and disclosing changes.

\section{Next experiments}
Freeze the current implementation. Test frozen native, proxy, and R0 interventions;
compare bounded-distance source maps without forcing an infeasible G16; obtain
quality--recovery Pareto curves; and confirm a small number of predeclared operating
points on unseen queries and matched cold/warm timing blocks. Do not hide the
current negative results, retune on the historical benchmark, or reinterpret failed
non-inferiority tests as passes.

\bibliographystyle{plain}
\bibliography{references}
\end{document}
